%%% FP2001.tex 05.10.2000 - Vesa Kolhinen <vesa.kolhinen@phys.jyu.fi>
%%% Original version:
%%% FP2000.tex 12.10.1999 - Eero Noponen <eero.noponen@hut.fi>
%%%
%%%
%%%
%%% Sample LaTeX file for the abstracts for the Proceedings of the
%%% Annual Meeting of the Finnish Physical Society.

\documentclass[12pt]{article}
\usepackage{fp2009}
\usepackage[T1]{fontenc}
\usepackage[latin1]{inputenc}
\usepackage{hyperref}

%%% Remove the following line if Times font is unavailable.
\usepackage{times}

%%% The following line is required if the document contains figures.
\usepackage{graphicx}

%%% The following line is for LaTeX 2.09.
%\documentstyle[FP2001,12pt]{article}

\begin{document}

%%%
%%% The title of the abstract is given here as \title{TITLE TEXT}.
%%%
\title{HOW TO WRITE ABSTRACTS FOR THE PROCEEDINGS OF THE XXXIX ANNUAL CONFERENCE
OF THE FINNISH PHYSICAL SOCIETY}

%%%
%%% The list of authors; underline the name of the author who is responsible
%%% for the presentation of the paper.
%%%
\underline{A. Fyysikko} and B. Fyysikko

%%%
%%% The address.
%%%
Finnish Physical Society, P.O.B. 64, FIN-00014 University of Helsinki, Finland\\
email: toimisto@fyysikkoseura.fi

%%% Do not remove.
\vspace{\baselineskip}

%%%
%%% Body text begins here.
%%%

Abstracts will be published in the \emph{Proceedings of the XLIII
Annual Conference of the Finnish Physical Society}, Report Series of
the Helsinki University of Technology. The language of the abstracts
is English.

Abstracts must be submitted to the organizers electronically, using
the file upload form at
\url{http://ltl.tkk.fi/wiki/Physics_Days_2009/Abstracts}. Files
prepared with \LaTeX\ in eps format or with Word in doc format are
accepted, as are also pdf files. Pdf format compiled using this
\LaTeX\ template are strongly preferred. Submitted files will be
published as they are. To maintain a coherent layout of the
proceedings all authors have to write the abstracts using this
sample file as a template. The sample file is designed to produce
correctly formatted abstracts on an A4 page with line spacing 1 and
3 cm wide margins. The font should be Times (New) Roman 12. Those
who use \LaTeX\ should prepare the final pdf file with the commands
\mbox{\tt dvips -o file.ps file.dvi} and \mbox{\tt ps2pdf file.ps}.
In the output, check that the possible web links are clickable and
lead to the correct address.

The title has to be written in bold capitals (the PACS code used
previously should not be included anymore). Title, authors, and
address information have to be separated from each other by an empty
line. The name of the author who will give the presentation has to
be underlined. The email address of this author can be given on a
separate line below the address information.

The text has to be justified. Formulas have to be centered and, if
numbering is necessary, the number has to be placed on the right in
parentheses. References are cited in the text like this \cite{ref1}.
The reference list should use the format shown below. You can add
web links using the example below --- if possible, please do this
for each reference.

%%% The following minipage structures are needed to position the figure
%%% beside the text paragraph.
%%% For a standard figure environment, see the example below.

\begin{minipage}{9cm}
Figures may be included in the submitted files. Here we show how a
figure can be included in a \LaTeX\ file.
\end{minipage}
\hfill
\begin{minipage}{5cm}
\includegraphics[width=4cm]{logo.eps}
\end{minipage}

%%%
%%% A more typical Figure.
%%%
%\begin{figure}[!h]
%\centering\includegraphics[width=4cm]{logo.eps}
%\caption{Mallikuva.}
%\end{figure}

%%%
%%% The bibliography: the references are listed here.
%%%
\begin{thebibliography}{9}

\bibitem{ref1}
A. Fyysikko and B. Fyysikko,
\href{http://link.aps.org/abstract/PR/v108/p1175}{Fysiikan Lehti 108
(1957) 1175}.
\end{thebibliography}

\end{document}
